% ============================================================
%  Balancing Robot – API & Benutzerhandbuch
%  Simples, stabiles Design. Compiler: pdflatex
% ============================================================
\documentclass[11pt, a4paper]{scrreprt}

\usepackage[T1]{fontenc}
\usepackage[utf8]{inputenc}
\usepackage[ngerman]{babel}
\usepackage[a4paper, top=25mm, bottom=30mm, left=25mm, right=25mm]{geometry}
\usepackage{xcolor}
\usepackage{graphicx}
\usepackage{helvet}
\usepackage{booktabs}
\usepackage{enumitem}
\usepackage{listings}
\usepackage{hyperref}
\usepackage[automark]{scrlayer-scrpage}

\renewcommand{\familydefault}{\sfdefault}

\definecolor{htbblue}{HTML}{0078C8}
\definecolor{htbgray}{HTML}{F4F4F4}

% Kopfzeile
\clearpairofpagestyles
\ohead{\pagemark}
\ihead{\headmark}
\KOMAoptions{headsepline=true}

% Überschriften blau
\addtokomafont{chapter}{\color{htbblue}}
\addtokomafont{section}{\color{htbblue}}
\addtokomafont{subsection}{\color{htbblue}}

\hypersetup{colorlinks=true, linkcolor=htbblue, urlcolor=htbblue}

% Code-Listings
\lstset{
  language=Python,
  basicstyle=\ttfamily\small,
  keywordstyle=\color{htbblue}\bfseries,
  backgroundcolor=\color{htbgray},
  frame=single, framerule=0pt,
  numbers=left, numberstyle=\tiny\color{gray},
  breaklines=true,
  showstringspaces=false,
}

% ============================================================
\begin{document}

% ── Titelseite ───────────────────────────────────────────────
\begin{titlepage}
  \vspace*{40mm}
  {\color{htbblue}\rule{\textwidth}{1.5pt}}
  \vspace{6mm}

  {\fontsize{32}{38}\selectfont\bfseries Balancing Robot\par}
  \vspace{4mm}
  {\Large\color{htbblue} API \& Benutzerhandbuch\par}
  \vspace{6mm}
  {\color{htbblue}\rule{\textwidth}{0.5pt}}
  \vspace{8mm}

  {\normalsize
    HTW Saar – Industrielle Robotik · WiSe 2025/26\\[4pt]
    \textbf{Version:} 1.0\quad\textbf{Datum:} \today
  }

  \vfill
  \IfFileExists{htw_logo.png}{%
    \includegraphics[height=16mm]{htw_logo.png}%
  }{}
\end{titlepage}

% ── Inhaltsverzeichnis ───────────────────────────────────────
\tableofcontents

% ============================================================
\chapter{Schnellstart}

\section{Voraussetzungen}
% TODO

\section{Laptop-Test ohne Hardware}
% TODO

\section{Start auf dem Raspberry Pi}
% TODO

% ============================================================
\chapter{Modul: balancing\_platform}

\section{Klasse Platform}
% TODO

\section{Methoden}
% TODO

% ============================================================
\chapter{Modul: hardware}

\section{MockServoDriver}
% TODO

\section{PCA9685Driver}
% TODO

% ============================================================
\chapter{Modul: kinematik}

\section{KinematicsEngine}
% TODO

\section{KinematicsResult}
% TODO

\section{PlatformGeometry}
% TODO

% ============================================================
\chapter{Modul: vision}

\section{BallDetectorRuntime}
% TODO

\section{BallTracker}
% TODO

\section{HsvRange}
% TODO

% ============================================================
\chapter{Konfiguration}

\section{config.json}
% TODO

\section{positions.json}
% TODO

% ============================================================
\chapter{Inbetriebnahme}

\section{Installation}
% TODO

\section{Kamerakalibrierung}
% TODO

\section{Kalibrierung der Servos}
% TODO

\end{document}
